% Ad Atticum 1.12 (ad-atticum-01-12)
% Source: https://raw.githubusercontent.com/PerseusDL/canonical-latinLit/master/data/phi0474/phi057/phi0474.phi057.perseus-lat1.xml
% Fetched by scripts/fetch_latin.py

% Dateline (Perseus): Scr. Romae In Ian. a. 693 (61)

\ciceroLetterOpener{CICERO ATTICO salutem}

\ciceroSection{1} Teucris illa lentum sane negotium, neque Cornelius ad Terentiam postea rediit. opinor ad Considium, Axium, Selicium confugiendum est; nam a Caecilio propinqui minore centesimis nummum movere non possunt. sed ut ad prima illa redeam, nihil ego illa impudentius, astutius, lentius vidi. libertum mitto, Tito mandavi. σκήψεισ atque ἀναβολαί ; sed nescio an ταὐτόματον . nam mihi Pompeiani prodromi nuntiant aperte Pompeium acturum Antonio succedi, oportere eodemque tempore aget praetor ad populum. res eius modi est ut ego nec per bonorum nec per popularem existimationem honeste possim hominem defendere, nec mihi libeat, quod vel maximum est. etenim accidit hoc, quod totum cuius modi sit mando tibi ut perspicias.

\ciceroSection{2} libertum ego habeo sane nequam hominem, Hilarum dico, ratiocinatorem et clientem tuum. de eo mihi Valerius interpres nuntiat Thyillusque se audisse scribit haec, esse hominem cum Antonio; Antonium porro in cogendis pecuniis dictitare partem mihi quaeri et a me custodem communis quaestus libertum esse missum. non sum mediocriter commotus neque tamen credidi, sed certe aliquid sermonis fuit. totum investiga, cognosce, perspice et nebulonem illum, si quo pacto potes, ex istis locis amove. huius sermonis Valerius auctorem Cn. Plancium nominabat. mando tibi plane totum ut videas cuius modi sit.

\ciceroSection{3} Pompeium nobis amicissimum constat esse. divortium Muciae vehementer probatur. P. Clodium Appi f., credo te audisse cum veste muliebri deprehensum domi C. Caesaris cum pro populo fieret, eumque per manus servulae servatum et eductum; rem esse insigni infamia. quod te moleste ferre certo scio.

\ciceroSection{4} quid praeterea ad te scribam non habeo, et me hercule eram in scribendo conturbatior. nam puer festivus anagnostes noster Sositheus decesserat meque plus quam servi mors debere videbatur commoverat. tu velim saepe ad nos scribas. si rem nullam habebis, quod in buccam venerit scribito. Kal. Ianuariis M. Messalla, M. Pisone coss.
